\section{Global regularity changes the minimal lift resources}\label{sec:regularity}
The generic selectors used above are guaranteed to be differentiable only
on a dense open set.
Requiring $C^1$ factor maps on the entire sphere is a stronger
condition. We determine its exact cost in the number, total capacity,
and total dimension of Lorentz factors in the full-row
identities~\eqref{eq:full-rows}, including factors that vanish on open
sets.

\subsection{Lorentz phases and pointwise no-sharing}
We use the Euclidean self-dual pairing on
$Q_m=\{(t,z):t\geq\norm z\}$. A fixed rescaling converts it to the Jordan
trace pairing and changes no rank or resource count.
Suppose~\eqref{eq:full-rows} has globally $C^1$ factors in $Q_{m_i}$,
$m_i\geq3$, and put $r_i=m_i-2$, the capacity of factor $i$; we call
$i$ a \emph{label}. At a contact point define
\[
 H_i^a(x)(u,v)=-\ip{D_aX_i(x)[u]}{DY_i^a(x_a)[v]}.
\]
Then $\sum_iH_i^a$ is the sphere metric, of rank $p_a=s_a-1$.
If $X_i$ or $Y_i^a$ is zero at the contact, so is its derivative, by
pointedness, and $H_i^a=0$. If one factor is interior, the other is zero.
Otherwise, on a neighborhood where both factors remain nonzero,
complementarity gives
\[
 X_i=\alpha_i(1,q_i^a(x_a)),\qquad
 Y_i^a=\beta_i^a(1,-q_i^a),\qquad q_i^a\in\Sph^{r_i},
 \quad\alpha_i,\beta_i^a>0.
\]
Complementarity holds for every $x$ with the same $x_a$ (the cylinder
identity), so the primal phase $q_i^a$ depends only on $x_a$ there.
Differentiation cancels every term containing a derivative of the
amplitudes, and
\begin{equation}\label{eq:phase-gram}
 H_i^a=\alpha_i\beta_i^a(Dq_i^a)^*Dq_i^a\succeq0,
 \qquad\rank H_i^a\leq r_i.
\end{equation}
Label $i$ is \emph{active} for source $a$ at $x$ if $H_i^a(x)\ne0$.
Then the primal phase is locally constant in every other source
direction, so at each point a label is active for at most one source
(\emph{pointwise no-sharing}):
\begin{equation}\label{eq:point-no-sharing}
 H_i^a(x)\ne0\quad\Longrightarrow\quad H_i^b(x)=0\quad(b\ne a).
\end{equation}
A label may be active for different sources at different points.

\subsection{A covering argument with top cohomology classes}
The lower bounds use three standard topological facts, with real
cohomology coefficients. First, if open sets $U_a$ cover a
space $M$ and classes $u_a\in H^*(M)$ satisfy
$u_a|_{U_a}=0$, the relative cup product implies
$\prod_au_a=0$: lift each class to $H^*(M,U_a)$ and multiply to get a
class in $H^*(M,\bigcup_aU_a)=0$. The same argument applies to any
subfamily that covers. Second, by the K\"unneth formula, on a product of
spheres the product of any subset of the pulled-back top classes is
nonzero. Third, a proper open subset of a connected closed oriented $p$-manifold has
zero ordinary cohomology in degree $p$ when every component is
noncompact; in particular this holds for every proper open subset of
$\Sph^p$. See~\cite[Sections~3.2--3.3]{Hatcher2002}.

\begin{lemma}[Saturated labels have a proper phase domain]\label{lem:phase-domain}
Fix a source $a$ with $p_a\geq2$ and a compact set $E\subset M$ on which
the labels active for $a$ are exactly those in $J$, with $\sum_{i\in J}r_i=p_a$ and
$r_i<p_a$ for every $i\in J$. There is an open neighborhood $P$ of
$\pi_a(E)$ in $\Sph^{p_a}$ such that the dual phases give a local
$C^1$ diffeomorphism
\[
 (q_i^a)_{i\in J}:P\longrightarrow\prod_{i\in J}\Sph^{r_i}.
\]
Moreover, $P\ne\Sph^{p_a}$, so the pulled-back top class of the source
sphere vanishes on $\pi_a^{-1}(P)$.
\end{lemma}
\begin{proof}
At a point of $E$, every active primal and dual factor is nonzero.
Fix all other primal coordinates and vary $x_a$ nearby. By continuity the primal factor stays nonzero, so the dual factor stays
on the nonzero boundary. Hence its normalized phase is defined
on an open neighborhood of $x_a$. The equality $\sum_{i\in J}r_i=p_a$,
together with~\eqref{eq:phase-gram}, makes
the joint phase derivative invertible. Invertibility is open. The union
of these phase domains, restricted to their common invertibility locus,
provides $P$.

If $P$ were the whole sphere, compactness would make this local
diffeomorphism a finite covering of the connected target. If any
$r_i=1$, the target has infinite fundamental group, which is
incompatible with a finite covering by the simply connected $\Sph^{p_a}$.
Otherwise the target is simply connected, so the covering is a
diffeomorphism. But the target has at least two factors, since
$r_i<p_a$, so its intermediate cohomology differs from that of a sphere.
Hence $P$ is proper, and the top-cohomology fact above completes the
proof.
\end{proof}

\begin{theorem}[Exact smooth Lorentz resources]\label{thm:smooth-lorentz}
Fix an integer $c\geq1$ and allow Lorentz dimensions $3\leq m_i\leq c+2$.
For full-row factorizations~\eqref{eq:full-rows} with globally $C^1$
primal and dual maps, define
\[
 h_a=\begin{cases}1,&p_a\leq c,\\
 \lceil(p_a+1)/c\rceil,&p_a>c,\end{cases}
 \qquad
 \tau_a=\begin{cases}p_a,&p_a\leq c,\\p_a+1,&p_a>c.\end{cases}
\]
Let $L$ be the number of non-ray factors, $T=\sum_i(m_i-2)$ their total
capacity, and $D=\sum_im_i$ their total ambient dimension. Their exact
minima, attained simultaneously, are
\begin{equation}\label{eq:smooth-frontier}
 L_{\min}=\sum_ah_a,\qquad T_{\min}=\sum_a\tau_a,
 \qquad D_{\min}=\sum_a(\tau_a+2h_a).
\end{equation}
The minima are attained by Slater affine lifts with genuine whole-slice
certificates. Additional ray factors do not lower them.
\end{theorem}
\begin{proof}
Write
$C_a=\sum_{i\text{ active for }a}r_i$, $N_a=\#\{i\text{ active for }a\}$.
The metric identity and~\eqref{eq:point-no-sharing} give
\[
 C_a\geq p_a,\quad N_a\geq\lceil p_a/c\rceil,
 \qquad\sum_aC_a\leq T,\quad\sum_aN_a\leq L.
\]
Rays have $H_i^a=0$ and are omitted.

Let $A=\{a:p_a>c\}$. If $T<\sum_a(p_a+\mathbf1_{a\in A})$, at every
point some $a\in A$ has $C_a=p_a$. Thus the sets
$E_a=\{C_a=p_a\}$, $a\in A$, cover $M$. They are closed: each active
locus is open, and $C_a\geq p_a$. Partition $E_a$ by its exact set of
active labels. This is a finite clopen partition into compact sets
$E_{a,J}$: active labels remain active nearby, and an additional active
label would increase $C_a$. Lemma~\ref{lem:phase-domain} supplies a
proper phase neighborhood $P_{a,J}$ for each compact piece. Choose
pairwise disjoint open neighborhoods $W_{a,J}$ of the pieces, small
enough to lie in $\pi_a^{-1}(P_{a,J})$. The source top class vanishes
on $U_a=\coprod_JW_{a,J}$. These sets cover $M$, so their relative cup
product vanishes, contradicting K\"unneth. Hence $T\geq\sum_a\tau_a$.

The pointwise count already gives $N_a\geq h_a$ except for
$A_0=\{a:p_a>c,\ c\mid p_a\}$. If $L<\sum_ah_a$, at every point
some $a\in A_0$ has exactly $p_a/c$ active labels. These closed sets
cover $M$. Partition them by their active labels; on each piece every
label must
have capacity $c$, and their capacities sum to $p_a$. The same phase
and relative cup argument gives a contradiction. Thus $L\geq\sum_ah_a$.
Since $D=T+2L$, this also proves the dimension lower bound.

For attainment, if $p_a\leq c$ use the direct trace section
$X_a=(1,x_a)\in Q_{p_a+2}$, with certificate $(1,-v)$ in its row
and zero in other rows. Otherwise partition all $s_a=p_a+1$ coordinates
into $h_a$ nonempty groups of size at most $c$, and impose
\begin{equation}\label{eq:grouped-lorentz}
 A_G(t_G,x_G)=\left(\frac{1+t_G}{2},\frac{1-t_G}{2},x_G\right)
 \in Q_{|G|+2},\qquad\sum_Gt_G=1.
\end{equation}
This is equivalent to $t_G\geq\norm{x_G}^2$, hence projects exactly
onto the ball and is Slater at $x=0$ with all $t_G>0$. On the boundary
choose $t_G=\norm{x_G}^2$. The polynomial certificate
\[
 B_G(v)=\left(\frac{1+\norm{v_G}^2}{2},
 -\frac{1-\norm{v_G}^2}{2},-v_G\right)
\]
has pairing $(t_G+\norm{v_G}^2)/2-\ip{x_G}{v_G}$.
Summing gives $1-\ip xv$ on the entire slice. Each group contributes
one factor and capacity $|G|$. Separate factors for different sources
attain all three claimed minima.
\end{proof}

For comparison, among finite exact relative-Slater affine lifts without
a global $C^1$ requirement, generic definable selectors and the
pointwise argument give the exact minima
\begin{equation}\label{eq:nonsmooth-product}
 L_0=\sum_a\lceil p_a/c\rceil,\qquad T_0=\sum_ap_a,
 \qquad D_0=T_0+2L_0.
\end{equation}
Independent norm trees attain them by the construction in the proof of
Corollary~\ref{cor:dimension-resources}: partition $p_a$ into $\lceil p_a/c\rceil$
positive integers at most $c$; a rooted tree with one internal node with
$k+1$ children for each part $k$ has $1+p_a=s_a$ leaves. Pointwise
no-sharing at a common generic regular contact proves the matching lower
bounds. Such a contact
lies in the primal regular locus and in the inverse images, under each
source projection, of the dual regular loci; by definable selection this
finite intersection is dense and open. The comparison makes no claim
about arbitrary factor maps. Thus global regularity adds one capacity
unit for each source with $p_a>c$, and one factor precisely for those
sources with $p_a>c$ and $c\mid p_a$. These are formulation resources,
not intrinsic barrier parameters of the projected body: the standard
parameter $2L$ of a product of $L$ Lorentz factors can decrease on
the affine slice, as examined below.

\subsection{Why local and joint identities are weaker}
If only a positive weighted sum of the full labelled rows must factor,
fewer factors can suffice.
\begin{proposition}[Joint weighted rows have a different sharp count]\label{prop:joint}
For two balls of common dimension $s\geq3$ and positive weights
$\lambda_1+\lambda_2=1$, the kernel
\[
 S_\lambda(x,z)=\sum_{a=1}^2\lambda_a(1-\ip{x_a}{z_a})
\]
on $(\Sph^{s-1})^2$ has globally $C^1$ factorizations over copies of
$Q_3$, and the least number of factors in them is $2s-1$. The full
labelled rows~\eqref{eq:full-rows} require $2s$ factors.
\end{proposition}
\begin{proof}
Mixed differentiation of $S_\lambda$ at a contact gives a metric
$g_\lambda=\sum_iH_i$ of rank $n=2(s-1)$, so $L\geq n$.
If $L=n$, every $H_i$ has rank one everywhere and both factors of each
label are nonzero. The phases $q_i:M\to\Sph^1$ lift to $C^1$
functions $\theta_i:M\to\R$ since $M$ is simply connected. The identity
$g_\lambda=\sum_i\alpha_i\beta_i\,d\theta_i^2$ makes
$(\theta_1,\ldots,\theta_n):M\to\R^n$ a local diffeomorphism.
Its image would be both nonempty open and compact, which is impossible.
Thus $L\geq n+1=2s-1$.

For the upper bound put
$Z(x)=(\sqrt{\lambda_1}x_1,\sqrt{\lambda_2}x_2)\in\Sph^{2s-1}$.
Choose a unit vector $q$ supported in the first coordinate block, and set
\[
 c(x)=1-\ip q{Z(x)}>0,\qquad
 f(x)=\frac{Z(x)-\ip q{Z(x)}q}{c(x)}\in q^\perp\cong\R^{2s-1}.
\]
Direct expansion of squared norms gives the stereographic identity
\[
 1-\ip{Z(x)}{Z(z)}=\frac{c(x)c(z)}2\norm{f(x)-f(z)}^2.
\]
For $U(t)=((1+t^2)/2,(1-t^2)/2,t)$ and
$V(t)=((1+t^2)/2,-(1-t^2)/2,-t)$, both vectors are in $Q_3$ and
$\ip{U(t)}{V(u)}=(t-u)^2/2$. Thus
$c(x)U(f_j(x))$, $c(z)V(f_j(z))$, $1\leq j\leq2s-1$, give the
required smooth factors. The full-row count follows from
Theorem~\ref{thm:smooth-lorentz} with $c=1$.
\end{proof}

For $k\geq2$ equal balls, the same lower-bound argument gives
$k(s-1)+1$ factors. Pairing the sources and using the preceding
construction on each pair, with one coordinate-square construction on
any remaining source, gives $ks-\lfloor k/2\rfloor$ factors. We do not
determine the exact count for $k\geq3$. Both bounds concern only the
contact kernel; neither construction supplies the full labelled rows of
an affine lift. The phase argument alone cannot close the gap, because a
product of spheres of total dimension $n$ embeds in $\R^{n+1}$. To see
this, embed a sphere as an orientable hypersurface,
then observe that the boundary of a small tubular neighborhood of
$N^m\subset\R^{m+1}\subset\R^{m+p+1}$ is $N\times\Sph^p$;
the added normal bundle is trivial. Induction proves the assertion.

\begin{remark}[Switching must pass through vertex sets with interior]\label{rem:switching}
Call a label \emph{productive} for a source if it is active for that
source at some point. A Lorentz label productive for two sources must
have a primal vertex set $\{x:X_i(x)=0\}$ with nonempty interior. Indeed, choose dual phase patches for both sources
on which their derivatives are nonzero. Their images each contain more
than one point, so smaller patches can be chosen with disjoint images.
On the product of these patches, cylindrical complementarity would force
a nonzero primal phase to equal both images. Thus the primal factor is
zero on that open cylinder. Moreover, the nonzero locus of each
productive dual factor cannot be dense. Otherwise, the identity saying
that the primal derivative in every unrelated source direction is radial
extends by continuity to all points; this excludes productivity for the
second source. In particular, if the factor maps are analytic, or every
primal vertex set has empty interior, no label is productive for two
sources. The proof of
Theorem~\ref{thm:smooth-lorentz} does not need either extra hypothesis.
\end{remark}

\begin{example}[A codimension-two singular set removes the extra cost of smoothness]\label{ex:tree-regularity}
In a full binary norm tree for $\B_2^s$, let $t_v(x)=\norm{x_{S_v}}$
for each internal node $v$, where $S_v$ is its set of descendant leaves;
leaves carry their signed coordinates and the root value is one.
At each internal node with children $v_0,v_1$ take
\[
 A_v(x)=(t_v(x),t_{v_0}(x),t_{v_1}(x)),\qquad
 B_v(z)=(t_v(z),-t_{v_0}(z),-t_{v_1}(z)).
\]
Their pairings telescope to $1-\ip xz$. The dual factors are genuine
certificates of the tree slice, since the same telescoping works when
the primal internal variables are free. Boundary fibers are unique:
the nonnegative squared cone slacks sum to $1-\norm x^2=0$, forcing
every internal norm to its minimum. Both factor maps are globally
Lipschitz and semialgebraic, and analytic away from the sets
$x_{S_v}=0$. Every nonroot internal node has at least two leaves, so
these sets have codimension at least two in the boundary sphere when
nonempty. Nevertheless, the tree uses only $s-1$ copies of $Q_3$, whereas
for $s\geq3$ Theorem~\ref{thm:smooth-lorentz} requires $s$ copies for
globally $C^1$ factors. Thus unique boundary fibers, Lipschitz
regularity, and analyticity on a dense open set do not imply the
conclusion of Theorem~\ref{thm:smooth-lorentz}.
\end{example}

The tree also has a unique genuine dual certificate for every support
$z\in\Sph^{s-1}$. Write an arbitrary dual block at node $v$
as $(\alpha_v,\beta_{v,0},\beta_{v,1})\in Q_3$. Matching coefficients on
the entire affine slice forces
\[
 \alpha_{\mathrm{root}}=1,\qquad
 \beta_{v,j}=\begin{cases}
 -z_\ell,&\text{if child }v_j\text{ is leaf }\ell,\\
 -\alpha_{v_j},&\text{if child }v_j\text{ is internal}.
 \end{cases}
\]
Induction from the leaves, using
$\alpha_v\geq(\beta_{v,0}^2+\beta_{v,1}^2)^{1/2}$, gives
$\alpha_v\geq\norm{z_{S_v}}$. At the root this is equality because
$\norm z=1$. Equality then propagates down the tree: the head
(first coordinate) of each child block is nonnegative and at least its
subtree norm, so equality in the sum
of their squares forces equality for each child separately. This also
handles a zero subtree, whose head and descendant coordinates must all
vanish. Thus the unique tuple is exactly $B_v(z)$ above. Every nonzero
block has Jordan rank one; a zero block has rank zero.

\begin{corollary}[Exact worst pure-row certificate rank over real PSD2]
\label{cor:psd2-worst-pure-row}
For $C=\prod_{a=1}^h\B_2^{s_a}$, $s_a\geq2$, fix any positive weights
$\lambda_1,\ldots,\lambda_h$. Let $\cD_{\cL,a}(u_a)$ denote the full
fiber of genuine certificates of $1-\ip{x_a}{u_a}$ on the affine slice
of a lift $\cL$. Among finite exact relative-Slater lifts over real
$2\times2$ positive semidefinite cones and optional rays,
\begin{equation}\label{eq:psd2-worst-pure-row}
 \inf_{\cL}\ \sup_{u\in M}\;
 \sup_{\substack{Y^a\in\cD_{\cL,a}(u_a)\\1\leq a\leq h}}
 \sum_i\rank\left(\sum_a\lambda_aY_i^a\right)
 =\sum_a(s_a-1).
\end{equation}
\end{corollary}
\begin{proof}
The lower bound is Theorem~\ref{thm:product-existence} with $B=1$.
For the upper bound take independent binary norm trees, with $s_a-1$
factors for source $a$. The linear isomorphism
\[
 (t,z_1,z_2)\longmapsto
 \begin{pmatrix}t+z_1&z_2\\z_2&t-z_1\end{pmatrix}
\]
identifies $Q_3$ with the real PSD2 cone; multiplying the dual image by
$1/2$ preserves the pairing and does not change rank.

Any genuine pure-row certificate for source $a$ vanishes on all factors
belonging to other sources. Indeed evaluate its identity at a feasible
tuple with source $a$ at its support contact and every other source at
a Slater point of its tree. The slack is zero, and each nonnegative
block pairing must therefore vanish. A positive dual element orthogonal
to an interior primal element is zero. What remains on source $a$ is
its unique tree certificate, by the coefficient argument above.
Consequently every positive pure-row aggregate has rank at most
$\sum_a(s_a-1)$. Equality holds on the dense open support set where
every internal subtree norm is positive. This proves the upper bound
and hence the formula.
\end{proof}

The inner supremum in~\eqref{eq:psd2-worst-pure-row} ranges only over
genuine \emph{pure-row} certificates; the corollary neither minimizes
over nor describes all certificates of the aggregate objective.

Local saturation alone gives an exact linear-algebra statement. If positive semidefinite contact forms $H_i$ sum to a positive
metric $H$ and $\sum_i\rank H_i=\rank H$, then
$P_i=H^{-1/2}H_iH^{-1/2}$ are mutually orthogonal projections. Indeed
write $P_i=Z_iZ_i^*$ with $Z_i$ of full column rank; concatenating the
$Z_i$ gives a square matrix $Z$ with $ZZ^*=I$, hence $Z^*Z=I$.
This decomposes the tangent metric at the contact orthogonally but
says nothing about unused parts of the cone, higher
derivatives, or the global affine slice. The global arguments above use
compactness and the full two-point identities in addition to this local
rank equality.

\begin{proposition}[Everywhere differentiability already forbids saturation]\label{prop:differentiable}
Assume the Lorentz full-row factors are everywhere Fr\'echet
differentiable, without continuity of their derivatives, and $p_a\geq2$.
If $r_i\leq c<\max_ap_a$ for an integer $c\geq1$, then
\[
 T\geq\sum_ap_a+1.
\]
For one source this gives the same exact resource minima as
Theorem~\ref{thm:smooth-lorentz}. For several sources it gives the
additional count $L\geq\sum_a\lceil p_a/c\rceil+1$ whenever every
$p_a$ is divisible by $c$. We do not claim that, under differentiability
alone, the extra cost is additive over sources.
\end{proposition}
\begin{proof}
All pointwise phase and no-sharing calculations remain valid with
pointwise derivatives. Suppose $T=\sum_ap_a=:n$. Sum the row slacks
and aggregate their dual factors. At every contact, equality in the
rank bound forces every factor to be nonzero and to use its full
capacity. Its primal phase is therefore defined everywhere. Their
product $q:M\to\prod_i\Sph^{r_i}$ is differentiable with invertible
derivative, because its kernel is annihilated by the positive contact
metric. The finite-dimensional local inversion theorem for everywhere
differentiable maps with everywhere nonsingular derivative
\cite{SaintRaymond2002} makes $q$ a local homeomorphism. Compactness
makes it a finite covering. Since $M$ is simply connected, circle
factors are impossible; with all $r_i\geq2$ the covering is a
homeomorphism. Consequently
$\prod_a(1+t^{p_a})=\prod_i(1+t^{r_i})$.
The smallest exponent $j$ and its coefficient determine the
multiplicity of $1+t^j$; cancel that power and repeat. Thus the multisets of exponents
agree, contradicting the cap. This proves the strict capacity bound.
If every $p_a$ is divisible by $c$, the baseline count $L_0$ satisfies
$cL_0=\sum_ap_a$, so $T\leq cL$ forces $L\geq L_0+1$.
For one source, combine the capacity and count bounds with the grouped
polynomial construction. If $c\nmid p$, its factor count is already
$\lceil(p+1)/c\rceil$ by rounding; if $c\mid p$, use the extra factor.
\end{proof}

A smooth example shows why the proofs use the cover argument instead of
analytic continuation. Put $h(t)=e^{-1/t^2}$ for $t>0$ and $h(t)=0$ for
$t\leq0$, and use $U$ from Proposition~\ref{prop:joint}. For distinct
sources $a,b$, the factors
\[
 A(x)=h((x_a)_1)U((x_b)_1),\qquad
 B^a(v)=h(-v_1)U(0)
\]
are smooth and nonnegative Lorentz-valued, with
$\ip{A(x)}{B^a(x_a)}=0$ on the entire cylinder. The dual is not
identically zero, yet the primal phase varies in source $b$ wherever
$(x_a)_1>0$. Hence the identity that the primal derivative is radial
in other source directions, which holds where the dual factor is nonzero, cannot in general be extended across an open
set where the dual factor vanishes. The example refutes only this proof
step; it is not a full-row factorization with too few factors.
